Work in the \(r\) free coordinates of \(\mathcal G_{n,l}\). Let \(h^0\) equal one on every coordinate of arm \(1\) and zero on every free coordinate of arms \(2,\ldots,l\). Every type row selects exactly one observed-cell coordinate in each arm. The arm-one selected coordinate is free, whereas \(h^0\) vanishes on all later-arm free coordinates. Hence
\[(h^0)^{\mathsf T}\alpha_\tau=1\qquad\text{for every }\tau\in\mathcal T_{n,l}.\tag{1}\]
Because \(\mu\) is a vertex of the full-dimensional gauge section, its active row normals span \(\mathbb R^r\). Thus \(C_\mu\) is full-dimensional. If \(x=\sum_\tau q_\tau\alpha_\tau\in C_\mu\) with \(q_\tau\ge0\), then (1) gives
\[(h^0)^{\mathsf T}x=\sum_\tau q_\tau.\tag{2}\]
Therefore \((h^0)^{\mathsf T}x>0\) for every nonzero \(x\in C_\mu\), so \(C_\mu\) is pointed. Equation (2) also proves
\[C_\mu\cap\{x:(h^0)^{\mathsf T}x=1\}=\mathcal A_\mu.\tag{3}\]
This is a compact cross-section because it is the convex hull of finitely many rows.

The tangent cone is
\[T_\mu=\{h:\alpha_\tau^{\mathsf T}h\ge0\text{ for every }\tau\in\operatorname{Sig}(\mu)\}=C_\mu^*.\tag{4}\]
For nonzero \(h\in T_\mu\), put \(F_h=C_\mu\cap\{x:h^{\mathsf T}x=0\}\). The ray \(\mathbb R_+h\) is extreme in \(T_\mu\) if and only if the rows satisfying \(\alpha_\tau^{\mathsf T}h=0\) have rank \(r-1\). Indeed, rank \(r-1\) makes their common nullspace one-dimensional, whereas smaller rank supplies a nonzero direction transverse to \(h\) along which sufficiently small perturbations in both signs remain in the finitely described cone. The same rank is \(\dim F_h\), so this condition is equivalent to \(F_h\) being a facet of \(C_\mu\). Intersecting with the cross-section (3) gives a facet of \(\mathcal A_\mu\), and coning a facet of \(\mathcal A_\mu\) reverses the argument. This proves the facet--ray bijection. A ray of \(T_\mu\) is an incident bounded-edge direction exactly when its line search in \(\mathcal H^0_{n,l}\) has a finite endpoint, so the finite tangent rays form the asserted subset.

It remains to prove the product-box description. For an outcome pair \((a,b)\), write \(\eta(0)=a\), \(\eta(1)=b\), and define
\[R_{ab,z}(\mu)=\operatorname*{argmin}_{d\in\{0,1\}}\mu_{(\eta(d),d,z)}.\tag{5}\]
The separable-minimum identity in the global-parent theorem shows that the pair is active precisely when its minimum slack is zero, and then its active words are exactly \(\prod_zR_{ab,z}(\mu)\). Let \(e_{a,0,z}\) and \(e_{b,1,z}\) denote the corresponding coordinate rows after deletion of fixed gauge coordinates, with a deleted gauge coordinate represented by the zero row. For such an active pair, distributivity of convex hull over independent finite sums gives
\[\operatorname{conv}\left\{\sum_{z=1}^le_{\eta(\delta(z)),\delta(z),z}:\delta(z)\in R_{ab,z}(\mu)\right\}
=\sum_{z=1}^l\operatorname{conv}\left\{e_{\eta(d),d,z}:d\in R_{ab,z}(\mu)\right\}.\tag{6}\]
Every summand on the right is a singleton or a segment. The convex hull over the at most \(n^2\) active pairs is \(\mathcal A_\mu\). To optimize a rational linear objective, optimize its chosen coordinate independently in every arm of (6), add the \(l\) resulting values, and compare the answers across active pairs. The sets in (5), active-pair tests, and all rational comparisons require \(O(n^2l)\) operations on polynomial-length inputs. Thus exact linear optimization is polynomial, as claimed.